\section*{Capítulo \ref{ch:g12:randvar} --- Variables aleatorias y distribución binomial}

\begin{solution}{exo:g12:randvar:1}
$\E(X) = -0.3 + 0 + 0.8 + 0.5 = 1$.
$\E(X^2) = 0.3 \times 1 + 0 + 0.4\times4 + 0.1\times25 = 4.4$, luego por
König--Huygens $\V(X) = 4.4 - 1 = 3.4$ y
$\sigma(X) = \sqrt{3.4} \approx 1.84$.
\end{solution}

\begin{solution}{exo:g12:randvar:2}
\emph{1.} Las $10$ preguntas son \omterm{def:g12:randvar:bernoulli}{pruebas de Bernoulli} \omterm{def:g12:condprob:indep}{independientes} con
$p = \frac14$: $X \sim \mathcal B\left(10, \frac14\right)$,
$\E(X) = 2.5$ y
$\sigma(X) = \sqrt{10 \times \frac14 \times \frac34} = \sqrt{1.875}
\approx 1.37$.

\emph{2.} $\P(X = 0) = \left(\frac34\right)^{10} \approx 0.056$;
\[
\P(X = 5) = \binom{10}{5}\left(\frac14\right)^5\left(\frac34\right)^5
\approx 0.058;
\quad
\P(X \geq 1) = 1 - \left(\tfrac34\right)^{10} \approx 0.944 .
\]
\end{solution}

\begin{solution}{exo:g12:randvar:3}
Pruebas idénticas e \omterm{def:g12:condprob:indep}{independientes}: $X \sim \mathcal B(400,\ 0.05)$, luego
$\E(X) = 20$ partes y
\[
\sigma(X) = \sqrt{400 \times 0.05 \times 0.95} = \sqrt{19} \approx 4.4 .
\]
\end{solution}

\begin{solution}{exo:g12:randvar:4}
El pago $P$ recibido cumple $\E(P) = \frac{5 + 6}{6} = \frac{11}{6}$, así
que el precio equitativo es $c = \frac{11}{6} \approx 1.83$ euros. La
ganancia equitativa $G = P - \frac{11}6$ toma los valores
$-\frac{11}{6}, \frac{19}{6}, \frac{25}{6}$ con \omterm{def:g10:proba:distribution}{probabilidades}
$\frac46, \frac16, \frac16$:
\[
\V(G) = \E(G^2)
= \frac{4 \times 121 + 361 + 625}{6 \times 36} = \frac{1470}{216}
= \frac{245}{36} \approx 6.8,
\qquad \sigma(G) \approx 2.6 .
\]
Un juego equitativo tiene ganancia \emph{\omterm{def:g11:stat:mean}{media}} nula, pero sus resultados
siguen fluctuando: equitativo no es lo mismo que sin riesgo.
\end{solution}

\begin{solution}{exo:g12:randvar:5}
$X \sim \mathcal B(8,\ 0.7)$.
\[
\P(X = 6) = \binom86 (0.7)^6(0.3)^2 \approx 0.296 ;
\]
\[
\P(X \geq 6) = \P(6) + \P(7) + \P(8)
\approx 0.296 + 8(0.7)^7(0.3) + (0.7)^8
\approx 0.296 + 0.198 + 0.058 = 0.552 ;
\]
$\P(X \geq 1) = 1 - (0.3)^8 \approx 0.99993$. Moda: $(n+1)p = 6.3$, así que
el valor más probable es $k^* = 6$ (\cref{exo:g12:randvar:7}).
\end{solution}

\begin{solution}{exo:g12:randvar:6}
El número de pasajeros que se presentan es $X \sim \mathcal B(104,\ 0.9)$;
el vuelo tiene sobreventa efectiva cuando $X \geq 101$:
\[
\P(X \geq 101) = \sum_{k=101}^{104}\binom{104}{k}(0.9)^k(0.1)^{104-k}
\approx 0.006 .
\]
Vender un $4\,\%$ más de billetes que plazas provoca un incidente solo en
torno al $0.6\,\%$ de los vuelos: la economía que hay detrás de la
sobreventa.
\end{solution}

\begin{solution}{exo:g12:randvar:7}
\[
\frac{\P(X = k+1)}{\P(X = k)}
= \frac{\binom{n}{k+1}}{\binom nk}\cdot\frac{p}{1-p}
= \frac{n - k}{k + 1}\cdot\frac{p}{1-p},
\]
usando $\binom{n}{k+1} = \binom nk \frac{n-k}{k+1}$. Este cociente es
$\geq 1$ si y solo si $(n-k)p \geq (k+1)(1-p)$, si y solo si
$np - k p \geq k - kp + 1 - p$, si y solo si $k \leq (n+1)p - 1$. Así que
las \omterm{def:g10:proba:distribution}{probabilidades} crecen estrictamente mientras $k + 1 \leq (n+1)p$ y
decrecen después: el \omterm{def:g10:functions:extrema}{máximo} se alcanza en $k^* = \floor{(n+1)p}$
(compartido con $k^* - 1$ cuando $(n+1)p$ es entero).
\end{solution}

\begin{solution}{exo:g12:randvar:8}
\emph{1.} Que la primera cara salga en el lanzamiento $k$ significa $k-1$
cruces y después cara: \omterm{def:g10:proba:distribution}{probabilidad}
$\left(\frac12\right)^{k-1}\cdot\frac12 = 2^{-k}$, para
$1 \leq k \leq 10$. \omterm{def:g10:proba:distribution}{Probabilidad} total de ganar:
$\sum_{k=1}^{10} 2^{-k} = 1 - 2^{-10} = \frac{1023}{1024}$ (el juego solo
se pierde con diez cruces seguidas).

\emph{2.} Ganancia esperada:
\[
\sum_{k=1}^{10} 2^k \cdot 2^{-k} = \sum_{k=1}^{10} 1 = 10 \text{ euros}.
\]
Sin el tope, la suma $\sum_{k\geq1} 1$ diverge: la ganancia esperada es
infinita, aunque el juego casi siempre paga una cantidad pequeña; es la
famosa \emph{paradoja de San Petersburgo}, que muestra que la \omterm{def:g12:randvar:exp}{esperanza},
por sí sola, no mide el valor de un juego.
\end{solution}

\begin{solution}{pb:g12:randvar:1}
\textbf{1.} $X \sim \mathcal B\!\left(3, \frac16\right)$:
$\P(X = 0) = \frac{125}{216}$,
$\P(X = 1) = \frac{75}{216}$,
$\P(X = 2) = \frac{15}{216}$,
$\P(X = 3) = \frac{1}{216}$.

\textbf{2.} $\E(X) = 3 \times \frac16 = \frac12$;
$\V(X) = 3 \times \frac16 \times \frac56 = \frac{5}{12}$.

\textbf{3.} Pago esperado $2\,\E(X) = 1$ euro frente a una apuesta de $1$
euro: ganancia $0$; exactamente equitativo (una rareza).

\textbf{4.} Sin reposición, las extracciones son dependientes (las
posibilidades de la segunda carta dependen de la primera): no es binomial;
falla la casilla de la \omterm{def:g12:condprob:indep}{independencia} de la lista. Con reposición: $n = 5$
fijo, $p = \frac14$ constante y extracciones \omterm{def:g12:condprob:indep}{independientes}: binomial.

\textbf{5.} $\E(X) = 6$ y $\sigma(X) = \sqrt{4.2} \approx 2.05$.
$\E(S) = 5 \times 6 - 10 = 20$; $\sigma(S) = 5\sigma(X) \approx 10.25$.

\textbf{6.} Número fijo de $n = 105$ pruebas (los billetes), cada una con
dos resultados (se presenta o no) y la misma $p = 0.9$, supuestas
\omterm{def:g12:condprob:indep}{independientes}: binomial. La confesión: las familias pierden los vuelos
\emph{juntas} y las tormentas vacían aviones enteros; la \omterm{def:g12:condprob:indep}{independencia} es
el acto de fe del modelo, y las ausencias correlacionadas engordan las
colas más de lo que promete la binomial.

\textbf{7.} $\E(X) = 94.5$;
$\sigma(X) = \sqrt{105 \times 0.9 \times 0.1} = \sqrt{9.45} \approx 3.07$.

\textbf{8.} $\P(X \geq 101) = \sum_{k=101}^{105}
\binom{105}{k} 0.9^k\, 0.1^{105 - k}$.

\textbf{9.} $\approx 0.0167$: en torno a un vuelo de cada sesenta deja a
alguien en tierra.

\textbf{10.} $\E(\text{dejados en tierra}) \approx 0.023$ pasajeros por
vuelo: indemnización esperada $\approx 19$ euros, frente a $1\,000$ euros
de ingresos adicionales. Sobrevender cinco billetes es rentabilísimo; de
ahí que lo haga toda aerolínea.

\textbf{11.} De la tabla de \omterm{def:g10:proba:distribution}{probabilidades} de cola: $n = 105$:
$1.7\,\%$; $n = 106$: $4.0\,\%$; $n = 107$: $8.1\,\%$. La mayor venta que
cumple la \omterm{def:g11:scal:dot}{norma} es de $n = 106$ billetes.

\textbf{12.} $z = \frac{100.5 - 94.5}{3.07} \approx 1.95$: el umbral está
dos desviaciones típicas por encima de la \omterm{def:g11:stat:mean}{media}, y la regla aproximada de
la campana (``cola superior de $2\sigma$ $\approx 2.5\,\%$'') se queda
cerca del $1.7\,\%$ exacto; la curva suave que persigue a la binomial es
la protagonista de los capítulos que vienen.

\textbf{13.} $0.98^{20} + 20 \times 0.02 \times 0.98^{19} \approx 0.94$: el
lote honrado pasa el $94\,\%$ de las veces.

\textbf{14.} $0.9^{20} + 20 \times 0.1 \times 0.9^{19} \approx 0.39$: el
lote malo se cuela aun así el $39\,\%$ de las veces.

\textbf{15.} Riesgo del productor: que se rechace un lote bueno
($\approx 6\,\%$); riesgo del consumidor: que se acepte un lote malo
($\approx 39\,\%$). Una \omterm{def:g10:proba:sample}{muestra} mayor (con un umbral de aceptación
proporcional) reduce los dos, a costa de más inspección: la calidad tiene
una línea de presupuesto.

\textbf{16.} $\V\!\left(\frac Xn\right) = \frac{\V(X)}{n^2} =
\frac{np(1-p)}{n^2} = \frac{p(1-p)}{n}$: \omterm{def:g11:stat:variance}{desviación típica}
$\sqrt{\frac{p(1-p)}{n}}$. Cuadruplica la \omterm{def:g10:proba:sample}{muestra} y reducirás el ruido a la
mitad: la ley del $\frac{1}{\sqrt n}$ de los \omterm{def:g10:numbers:interval}{intervalos} de fluctuación del
año 10, ahora demostrada.

\textbf{17.} El tope limita el pago a $2^{10} = 1024$ euros, así que cada
una de las diez rondas aporta exactamente $1$ euro de \omterm{def:g12:randvar:exp}{esperanza}:
$\E = 10$. La \omterm{def:g12:randvar:exp}{esperanza} infinita del juego sin tope venía enteramente de
pagos astronómicamente raros y astronómicamente grandes; poner un tope es
confesar que ningún banco paga $2^{50}$ euros. Precio equitativo de una
papeleta del juego con tope: $10$ euros; y de repente ya no hay paradoja
que valga.

\textbf{18.} Premio esperado por papeleta:
$\frac{100 + 2 \times 50}{200} = 1$ euro frente a una papeleta de $2$
euros: margen del $50\,\%$, dieciocho veces el $2.7\,\%$ de la ruleta. La
rifa sobrevive porque sus jugadores están donando a sabiendas: la
``pérdida'' es justamente el objetivo.

\textbf{19.} Siniestros esperados:
$1\,000 \times 0.01 \times 10\,000 = 100\,000$; primas: $130\,000$:
beneficio esperado de $30\,000$ euros. \omterm{def:g11:stat:variance}{Desviación típica} del total de
siniestros:
$10\,000 \times \sqrt{1\,000 \times 0.01 \times 0.99} \approx 31\,500$
euros; un solo año malo corriente se come todo el beneficio esperado. De
ahí las reservas de capital, el reaseguro y carteras mucho mayores que mil
pólizas: las aseguradoras viven de la ley de los grandes números y guardan
reservas contra su lentitud.

\textbf{20.} Primero, la lista de comprobación ($n$ fijo, la misma $p$, la
\omterm{def:g12:condprob:indep}{independencia} \emph{afirmada}) o no hay binomial que valga. Después, navega
con $\E \pm \sigma$: las \omterm{def:g11:stat:mean}{medias} sitúan y las desviaciones avisan. Después,
pon precio a las colas: las denegaciones de embarque, los lotes malos y los
años ruinosos viven más allá de $2\sigma$. Advertencias: la \omterm{def:g12:condprob:indep}{independencia}
es una promesa del que modeliza, no del mundo; y la \omterm{def:g12:randvar:exp}{esperanza} ignora
justamente lo que te destruye. Las páginas que le faltan al reglamento
(con qué rapidez se asientan las frecuencias y qué forma toman las
fluctuaciones) son los dos capítulos siguientes.
\end{solution}
